% Generated from validated Special post-draft synthesis. Do not hand-edit.
\section*{Retrospective Signals}
\addcontentsline{toc}{section}{Retrospective Signals}
\smallskip
\noindent \textbf{Reasoningが実行能力の一層になった} R1、o3系、Claude 3.7、Gemini 2.5を横断すると、thinkingは単独カテゴリよりtool useやagent executionを支えるcapabilityとして見える。 \hfill{\footnotesize p.~\pageref{pkg:reasoning}}
\par
\smallskip
\noindent \textbf{Agentのaction surfaceが分化した} browser操作、web research、developer orchestration、cloud codingが別々の形で立ち上がり、model responseの外側が実用能力を左右し始めた。 \hfill{\footnotesize p.~\pageref{pkg:agents}}
\par
\smallskip
\noindent \textbf{Releaseだけではmodel familyを追えなくなった} research preview、API family、GA、deprecationが短い間隔で続き、availability lifecycle自体が技術運用上の論点になった。 \hfill{\footnotesize p.~\pageref{pkg:lifecycle}}
\par
\smallskip
\noindent \textbf{Open-weightとmultimodalityが同時に広がった} MoE、portable model、native multimodality、generative mediaが並行し、weights、hosted API、media workflowを別軸で比較する必要が生じた。 \hfill{\footnotesize p.~\pageref{pkg:open} / p.~\pageref{pkg:multimodal}}
\par
\smallskip
\noindent \textbf{Execution拡大がcontrol boundaryを可視化した} web search、safeguard、interpretabilityはまだ分散した萌芽だったが、外部toolへ接続するsystemにcontrolとobservabilityが必要になる方向を示した。 \hfill{\footnotesize p.~\pageref{pkg:control}}
\par
\medskip
\begin{claimboundary}[Retrospective scope]
本号は2025年前期を後日確認可能になった一次情報も用いて再構成するRetrospective Specialである。Coverage Periodと制作時点を同一視せず、vendor / project / author claimの境界は各記事で明示する。
\end{claimboundary}
\medskip
\tableofcontents
